%% Proofs for ED26-P07-monotonicity-0001.
\begin{annexitem}{annex:p7.monotonicity.type-order}
 {thm:p07-monotonicity-type-order}{The recursive type order}
\begin{proof}
Define a relation \(R\) recursively: equal base types are related; no
base is related to an arrow; and two arrows are related exactly when their
argument types are related in the reverse direction, their result types in
the forward direction, and their markings in the forward direction.
Induction on type height proves reflexivity. For transitivity, a chain
of comparable arrows gives reverse chains in their arguments, forward
chains in their results and a chain of markings. The induction hypothesis
and transitivity of the marking order give the required arrow comparison.
Thus \(R\) is a preorder satisfying the defining closure rule, so the least
preorder \(\preceq\) is included in \(R\).
Conversely, induction on height expresses each \(R\)-comparison by one
application of the closure rule to its smaller comparisons; hence
\(R\subseteq\preceq\).
The same induction and antisymmetry of markings prove antisymmetry.

For the auxiliary join, the diagonal case is immediate. In its arrow
clause the arguments are literally the same type, the result types are
below their recursively defined join, and both markings are below their
join. Apply the characterized arrow order. This proves the asserted upper
bound property wherever the operation is defined.

One finite-stage description, if desired, starts with the identity on
\emph{all} types and retains that identity at every stage while closing
under the arrow constructor. The union of these stages is \(R\).
Retention of the base identities is necessary when an arrow has a base
type as one of its components.
\end{proof}
\end{annexitem}

\begin{annexitem}{annex:p7.monotonicity.coercion-polarity}
 {thm:p07-monotonicity-coercion-polarity}{Monotonicity exposed by restriction}
\begin{proof}
We have \(\rho\preceq\sigma\), so
\((\sigma\xrightarrow+b)\preceq(\rho\xrightarrow+b)\).
The abstraction has intrinsic type \(\sigma\xrightarrow+b\), because
the occurrence of its binder \(f\) is in function position and is positive.
It can therefore be supplied to \(F\).
Its free \(x\)-occurrence lies in the unspecified argument position of
\(f:\sigma\); abstraction preserves that polarity, and the positive
argument position of \(F\) leaves it unspecified.
Neither the abstraction nor its enclosing application is a beta or eta
redex, so the term is normal.

For \(a\in D_b\), let \(H_a\) be the denotation of that abstraction at
\(\varphi_x^a\), coerced to \(D_{\rho\xrightarrow+b}\).
For \(g\in D_\rho\), the abstraction clause and (C) give
\[
 H_a(g)=\pi_{\rho,\sigma}(g)(a)=g(a);
\]
forgetting a marking is the same-function inclusion.
If \(a\leq_b a'\), monotonicity of every \(g\in D_\rho\) gives
\(H_a(g)\leq_b H_{a'}(g)\). Pointwise order gives
\(H_a\leq H_{a'}\). The fixed interpretation of \(F\) is monotone, so
\(\sem M_{\varphi_x^a}\leq_b\sem M_{\varphi_x^{a'}}\).

For a concrete nonvacuous instance, use the full structures over
\(D_b=\{0,1\}\) with \(0<1\), and interpret
\(F(H)=H(\operatorname{id}_{D_b})\). This evaluation functional is
monotone, and \(\sem M_{\varphi_x^a}=a\).
Thus the semantic conclusion is realized by an ordinary finite base
model, despite the unspecified raw occurrence.
\end{proof}
\end{annexitem}

\begin{annexitem}{annex:p7.monotonicity.polarity}
 {thm:p07-monotonicity-polarity}{The two polarity implications}
\begin{proof}
Simultaneous structural induction proves the two conclusions and the
coincidence property: interpretation depends only on free variables.
Variables give the identity map or an independent constant; constants
give independent constants. These observations also establish the
coincidence property at the leaves.

Consider \(M(N)\), where the intrinsic function type has marking \(m\).
Put \(F_a=\sem M_{\varphi_x^a}\) and
\(u_a=\pi_{\operatorname{tp}(N),\sigma}
                  (\sem N_{\varphi_x^a})\).
Coercion preserves every inequality supplied by induction.
Suppose first that all relevant occurrences are positive and \(a\leq b\).
The function occurrences give \(F_a\leq F_b\) pointwise.
If \(m=+\), induction gives \(u_a\leq u_b\), and
\[
 F_a(u_a)\leq F_a(u_b)\leq F_b(u_b).
\]
If \(m=-\), induction instead gives \(u_b\leq u_a\); antitonicity again
gives the first inequality in the same display.
If \(m=\cdot\), the argument can have no free \(x\), since every such
occurrence would have total polarity \(\cdot\). Coincidence gives
\(u_a=u_b\), and the pointwise comparison suffices.

If all relevant occurrences are negative, \(F_b\leq F_a\).
For \(m=+\), induction gives \(u_b\leq u_a\), so
\(F_b(u_b)\leq F_b(u_a)\leq F_a(u_a)\).
For \(m=-\), induction gives \(u_a\leq u_b\), and antitonicity yields
the same displayed comparison. For \(m=\cdot\), the argument is again
independent of \(x\). This exhausts application, including absence of \(x\)
from either subterm. Coincidence is preserved by the application clause.

At \(\lambda y.M\), if \(y=x\), the outside value of \(x\) is overwritten
and the denotation is independent of it. Otherwise choose \(y\ne x\).
For every \(d\in D_{\operatorname{tp}(y)}\), the two assignment changes
commute:
\((\varphi_x^a)_y^d=(\varphi_y^d)_x^a\).
Apply the induction hypothesis to the body at \(\varphi_y^d\), then use
pointwise order on the function domain. This gives both signs and
coincidence under abstraction.

For a full prestructure, the recursive interpretation of an abstraction
is a function in the unrestricted arrow domain. The result just proved
shows it is in the required monotone or antitone subdomain whenever the
intrinsic marking demands this. Thus the recursion is well typed.
\end{proof}
\end{annexitem}

\begin{annexitem}{annex:p7.monotonicity.substitution}
 {thm:p07-monotonicity-substitution}{Substitution with type and sign refinement}
\begin{proof}
First note that marking composition preserves refinement in each argument:
if \(m'\sqsubseteq m\) and \(n'\sqsubseteq n\), then
\(m'\circ n'\sqsubseteq m\circ n\). This follows by checking signs, with
\(\cdot\) the only possible proper upper marking.

Prove type refinement and occurrence refinement together by induction.
At a variable they are the definition of substitution and the identity
law \(+\circ n=n\). At a constant there are no occurrences and the type
is unchanged. In an application suppose the original head type is
\(\sigma\xrightarrow m\tau\). Its substituted type is
\(\sigma_s\xrightarrow{m_s}\tau_s\), where
\(\sigma\preceq\sigma_s\), \(m_s\sqsubseteq m\), and
\(\tau_s\preceq\tau\), by the arrow characterization and induction.
If the argument type is \(\rho\preceq\sigma\), its substituted type
\(\rho_s\) satisfies
\(\rho_s\preceq\rho\preceq\sigma\preceq\sigma_s\).
The application therefore remains well typed and its result refines
\(\tau\). Occurrences in the head use the induction hypothesis directly.
For occurrences in the argument, compose the refined head marking with
the refined argument occurrence. The marking observation in the first
paragraph proves the required bound.

At an abstraction rename the binder fresh, retaining its type.
If its old marking is a sign, all its bound-variable occurrences have that
sign and at least one exists. The substitution images contain no free
occurrence of the fresh binder, so substitution neither introduces a
different signed occurrence of it nor removes its original variable
leaves. Induction preserves that sign. If the old marking is \(\cdot\),
any new marking is below it. Together with the refined result type this
proves the abstraction type bound. Occurrences of other free variables
retain their body polarities. This proves the two syntactic assertions.

For (S), use structural induction, renaming binders away from relevant
images. The variable case is the definition of \(\varphi_s\); the
constant case follows from its assignment-independent interpretation.
In the application case retain the types just named. Its substituted
value, promoted to \(\tau\), is
\[
 \pi_{\tau_s,\tau}\!\left(
  \sem{M[s]}_\varphi
       \bigl(\pi_{\rho_s,\sigma_s}(\sem{N[s]}_\varphi)\bigr)\right).
\]
By composition, the argument coercion factors through
\(\rho\), \(\sigma\), and \(\sigma_s\).
Equation (C) then identifies this expression with the promoted function
\(\pi_{\sigma_s\xrightarrow{m_s}\tau_s,\sigma\xrightarrow m\tau}
   (\sem{M[s]}_\varphi)\)
applied to
\(\pi_{\rho,\sigma}
  \pi_{\rho_s,\rho}(\sem{N[s]}_\varphi)\).
The two induction hypotheses yield exactly
\(\sem{M(N)}_{\varphi_s}\).

For an abstraction with fresh binder \(y:\sigma\), compare functions at
an arbitrary \(a\in D_\sigma\). Equation (C) promotes only the result
value, and the induction hypothesis for the body applies.
Freshness ensures that composing the assignment with substitution and
then changing \(y\) to \(a\) has the same effect as first changing \(y\)
and using the substitution that fixes the binder. Both sides of (S)
therefore have equal values at every \(a\), so they are equal functions.

A beta redex has the type of its body. Apply the proved type refinement
to its one-variable substitution. For a tracked signed occurrence in the
body away from the binder, substitution preserves the sign by the
occurrence bound. For an occurrence in the argument, its old total sign
is the product of the abstraction's marking and its argument polarity.
If that product is a sign, both factors are signs; every free occurrence
of the binder in the body has the abstraction's sign. Substitution gives
each corresponding copy that same product. Erased occurrences make no
remaining assertion. This proves the stated subject-reduction conclusion.
\end{proof}
\end{annexitem}

\begin{annexitem}{annex:p7.monotonicity.conversion}
 {thm:p07-monotonicity-conversion}{Substitution stability without typed normal-form uniqueness}
\begin{proof}
Let \(x_1,\ldots,x_n\) list the free variables of \(M,N\).
Repeated abstraction congruence gives
\[
 \lambda x_1\cdots\lambda x_n.M
 \approx
 \lambda x_1\cdots\lambda x_n.N .
\]
Rename these binders to variables \(y_i\) fresh for the finitely many
images \(s(x_j)\) and for the original bodies. Application congruence
allows both sides to be applied to \(s(x_1),\ldots,s(x_n)\).
Each application is well typed: its argument type is below the type
assigned to the corresponding binder. Repeated beta conversion gives
\(M[s]\approx N[s]\), because the fresh names make these sequential
substitutions the specified simultaneous substitution. This proof does
not assume that an individual eta contraction remains an exact-domain
eta contraction after substitution.

For beta termination, erase markings in all types. Comparable types have the
same erased arrow tree, by induction on the type order. A marked term
becomes an ordinary simply typed term, and every beta contraction
erases to a genuine ordinary beta contraction. Ordinary strong normalization
therefore precludes an infinite sequence. This does not assert uniqueness
of the marked binder annotations left in an irreducible term.
\end{proof}
\end{annexitem}

\begin{annexitem}{annex:p7.monotonicity.instantiation}
 {thm:p07-monotonicity-instantiation}{Instantiation and functional comparison}
\begin{proof}
For the first assertion, Equiv gives comparisons in both directions between
each pair of convertible representatives, at the declared common type.
Two uses of Trans replace the endpoints.

For the second assertion list the finitely many free variables of \(M,N\)
as \(x_1,\ldots,x_n\). Repeated Func turns the premise into a comparison
of \(\lambda x_1\cdots\lambda x_n.M\) and
\(\lambda x_1\cdots\lambda x_n.N\).
At each level use a common arrow type, taking marking \(\cdot\) if the
two intrinsic markings differ. No freshness condition relative to the
premise is needed: the premise is universal.

Alpha-rename these binders to \(y_1,\ldots,y_n\) fresh for all the
finitely many images \(s(x_j)\) and for the free variables in the bodies.
Repeated Point applies the abstractions, in order, to
\(s(x_1),\ldots,s(x_n)\). Their types are below the required argument
types. Repeated beta conversion and the first assertion replace both
endpoints by \(M[s]\) and \(N[s]\). The fresh names ensure that these
sequential beta substitutions implement the intended simultaneous
substitution, even when some image contains an original \(x_j\).
If the premise itself has a derivation from \(\Gamma\), replace its leaf
in this finite derivation by that derivation. This proves closure of
conclusions under substitution without changing \(\Gamma\).

For the third assertion choose a fresh variable \(v:\sigma\) outside
\(M,N\). The hypothesis includes
\(M(v)\precsim_\tau N(v)\). Func yields a comparison of their abstractions
at \(\sigma\xrightarrow m\tau\); this comparison type is legitimate by
the hypotheses on the two original arrow types. Eta conversion and the
first assertion give the desired comparison of \(M,N\).
\end{proof}
\end{annexitem}


\begin{annexitem}{annex:p7.monotonicity.term-inclusions}
 {thm:p07-monotonicity-term-inclusions}{Order preservation between term domains}
\begin{proof}
Induct on type height. At a base comparison the two types coincide.
Suppose
\(\sigma\xrightarrow m\tau\preceq
  \sigma'\xrightarrow{m'}\tau'\).
A pointwise comparison over all of \(T_\sigma\) also holds over its
subset \(T_{\sigma'}\). By induction its result comparisons persist from
\(T_\tau\) to \(T_{\tau'}\). Extensionality in the definition of a term
structure therefore gives the head comparison at
\(\sigma'\xrightarrow{m'}\tau'\). These are all type-order cases.
\end{proof}
\end{annexitem}

\begin{annexitem}{annex:p7.monotonicity.soundness}
 {thm:p07-monotonicity-soundness}{Soundness with universal premises}
\begin{proof}
Fix a coherent structure satisfying \(\Gamma\) under every assignment.
Lemma~\ref{thm:p07-monotonicity-long-coherence} supplies property (E).
Induct on a finite derivation, proving that its conclusion holds under
every assignment. Premise leaves hold by hypothesis.
Ref and Trans use the domain preorder. Pres uses monotonicity and
composition of coercions. Equiv uses
property (E).

For Point, regard both head values at the common arrow comparison type.
Their order is pointwise, so it persists upon application to the promoted
argument value. Equation (C) identifies these evaluated values with the
promoted interpretations of the two actual applications.
For Mono, the head has a positively marked intrinsic type and is
monotone; promote the two argument values to that argument type and apply
it. Anti is the same argument with order reversed.
For Func, the induction hypothesis holds at every
\(\varphi_x^a\), not just at \(\varphi\).
The abstraction clause and pointwise order give the inequality of the two
functions; (C) promotes their result values when necessary.
This proves every rule case.
\end{proof}
\end{annexitem}

\begin{annexitem}{annex:p7.monotonicity.ordinary-long}
 {thm:p07-monotonicity-ordinary-long}{Ordinary beta-eta confluence and long forms}
\begin{proof}
The only typed normalization input needed initially is beta strong
normalization. We first establish the beta-eta Church--Rosser property on
ordinary unmarked terms, following the parallel construction in
\cite[Section 4.4]{Selinger2013LambdaNotesV2}.
Use the five auxiliary parallel clauses on
Sheet~\ref{sheet:p07-monotonicity}, writing \(P\mathrel{\mathsf{Par}}Q\)
in this proof. These clauses act on ordinary syntax, not on marked
binder annotations.

Induction on terms gives parallel reflexivity. Induction on a single
compatible beta or eta contraction embeds it in the parallel relation:
the root cases use the respective parallel contraction with reflexivity
on subterms; context cases use parallel application or abstraction.
Conversely each parallel step is a finite beta-eta reduction: first reduce
its subterms by induction, then perform its root contraction if its last
clause is beta or eta. The two reflexive transitive closures therefore
coincide.

Parallel substitution satisfies
\[
 P\mathrel{\mathsf{Par}}P',\quad Q\mathrel{\mathsf{Par}}Q'
 \quad\Longrightarrow\quad
 P[Q/x]\mathrel{\mathsf{Par}}P'[Q'/x].
\]
To prove this, freshen binders and induct on the parallel derivation for
\(P\). A variable equal to \(x\) uses \(Q\mathrel{\mathsf{Par}}Q'\);
another variable or constant uses reflexivity. Application and abstraction
use their induction hypotheses. At a beta clause, apply both hypotheses
and the parallel beta clause, then use the capture-avoiding composition
identity
\(R[Q'/x][S[Q'/x]/y]=R[S/y][Q'/x]\) when \(y\) is fresh for \(Q'\).
At an eta clause, freshness ensures that substituting \(Q'\) does not
introduce its bound variable into the contracted head, so the parallel eta
clause applies. This proves every case. Parallel reduction also cannot
introduce a new free variable, by the same clause induction.

Define the complete development \(P^*\) recursively. At a root beta
redex use the developed body with the developed argument substituted.
At a root eta redex \(\lambda x.Rx\), use \(R^*\).
Otherwise develop the immediate subterms and preserve their constructor.
Atoms remain fixed. We prove
\(P\mathrel{\mathsf{Par}}P'\Rightarrow
  P'\mathrel{\mathsf{Par}}P^*\)
by induction on the size of \(P\).
For an atom it is reflexivity.
For an application-congruence step \(RN\mathrel{\mathsf{Par}}R'N'\):
if \(R\) is not an abstraction, application of the two induction
hypotheses reaches \(R^*N^*\).
If \(R=\lambda x.Q\), an abstraction-preserving parallel step has
\(R'=\lambda x.Q'\); the parallel beta clause reaches \(Q^*[N^*/x]\).
The other possibility is an eta contraction
\(R=\lambda x.Sx\mathrel{\mathsf{Par}}S'=R'\), with
\(S\mathrel{\mathsf{Par}}S'\) and \(x\notin\FV(S)\).
The smaller term \(Sx\) gives
\(S'x\mathrel{\mathsf{Par}}(Sx)^*\).
Parallel substitution then gives
\(S'N'\mathrel{\mathsf{Par}}(Sx)^*[N^*/x]\), the required development.

For an abstraction-preserving step
\(\lambda x.Q\mathrel{\mathsf{Par}}\lambda x.Q'\):
if the source is not an eta redex, lift the body induction hypothesis.
Otherwise \(Q=Rx\) with \(x\notin\FV(R)\).
If its parallel step preserves application, \(Q'=R'x\), use the body
induction hypothesis \(R'\mathrel{\mathsf{Par}}R^*\) and parallel eta.
If it contracts beta, \(R=\lambda y.S\) and
\(Q'=S'[x/y]\) with \(S\mathrel{\mathsf{Par}}S'\).
Freshness gives \(\lambda x.S'[x/y]=_\alpha\lambda y.S'\).
Now \(R\mathrel{\mathsf{Par}}\lambda y.S'\); induction on the smaller
\(R\) gives \(\lambda y.S'\mathrel{\mathsf{Par}}R^*\).
These exhaust the application \(Rx\).
A root parallel beta step reaches the complete development by the
parallel substitution lemma and the two induction hypotheses.
A root parallel eta step reaches it by the induction hypothesis for its
head. This proves the development property.

Any two parallel successors consequently reduce in parallel to \(P^*\).
The diamond property extends to finite paths: induct first on the length
of one path to move a single competing step across it, and then on the
length of the second path. Since the transitive closures coincide,
beta-eta reduction is confluent. An induction on a conversion chain,
joining at each reversed step, gives a common reduct of convertible
terms. No normalization assumption is used for this confluence argument.

Now return to ordinary simply typed terms. Beta-normalize first, using
the ordinary beta strong-normalization theorem. An eta contraction in a
beta-normal term preserves beta normality: its contracted head is neutral,
since otherwise the source body already contains a beta redex; replacing
the eta abstraction by that neutral head creates no new beta redex in an
outer application. Eta contraction strictly reduces the number of
constructors. Hence a beta-normal term has a beta-eta short normal reduct.
Confluence makes this short reduct unique up to alpha conversion.

Define long expansion on beta-normal ordinary terms as follows.
For an abstraction, expand its body recursively.
For a neutral head \(hN_1\cdots N_k\), first expand the existing arguments
recursively. If its result type is
\(A_1\to\cdots\to A_r\to b\), introduce fresh
\(u_1:A_1,\ldots,u_r:A_r\), append their recursively expanded occurrences,
and abstract over these variables.
The recursion on existing subterms is structural; expansion of the newly
introduced variables is a separate recursion on their proper component
types. Both are finite. The output is beta-normal and long and is
beta-eta convertible to the input.

Long expansion is unchanged by eta contraction between beta-normal terms.
At a root \(\lambda x.Rx\), the expansion of the left side introduces the
same typed argument prefix and head spine as the expansion of \(R\);
fresh bound names are the only difference. The other cases are its
abstraction or neutral-argument contexts, where the recursive expansion
has exactly the corresponding subcall. Induction on that context proves
invariance. A term already long is fixed by expansion.
If two long terms are beta-eta convertible, reduce each by eta to its
unique common short normal form and apply this invariance. Both original
long terms are its same long expansion. This proves uniqueness, including
the uniqueness step left as an exercise after Proposition 8.7 of the
supporting source. Inert constants are simply additional neutral heads;
none of the arguments changes.
\end{proof}
\end{annexitem}

\begin{annexitem}{annex:p7.monotonicity.long-coherence}
 {thm:p07-monotonicity-long-coherence}{Type-directed elaboration and coherence}
\begin{proof}
Erase all markings in types and retain their base/arrow tree.
Comparable types have identical erasures, by the structural type-order
lemma. A marked term and every marked conversion step erase to ordinary
well-typed syntax and conversion. Ordinary long forms exist and are unique
by Lemma~\ref{thm:p07-monotonicity-ordinary-long}.

The comparison-directed construction follows that finite ordinary long
expansion, retaining its marked expected types.
Beta-normalize first; this is well typed by substitution refinement.
At an expected arrow
\(\tau=\sigma\xrightarrow m\rho\), a term of intrinsic type below
\(\tau\) has type
\(\delta\xrightarrow n\epsilon\), where
\(\sigma\preceq\delta\), \(\epsilon\preceq\rho\), and
\(n\sqsubseteq m\). A fresh \(u:\sigma\) can therefore be applied to
it. At an expected base, the neutral head determines each receiving
argument type. These are exactly the ordinary expansion calls with the
determined marked types restored; no arbitrary alternation of eta
expansion and contraction is performed.

We verify type validity together with signed-occurrence survival.
For a beta contraction, a tracked signed occurrence in the abstraction
body, distinct from its binder, survives capture-avoiding substitution.
For an occurrence in its argument, a signed total polarity requires
a signed abstraction marking and a signed argument polarity.
The binder of that abstraction is nonvacuous and all its occurrences
have the abstraction's sign. At least one copy therefore survives, and
each copy has the original product sign, by the substitution lemma.
The signed occurrences in a surrounding context also keep their signs:
a refined signed marking must be that same minimal marking.
An eta expansion retains all old free occurrences in its head and
introduces only a fresh bound variable.
Processing an existing abstraction at a smaller expected argument type
uses a fresh variable with that smaller type, precisely a type-refining
substitution. Its signed residuals are covered by the same argument.
These observations induct over every normalizing call.

In particular, the fresh \(u\) in the arrow case has sign \(n\) if
\(n\) is signed, and at least one occurrence survives.
The final abstraction marking is then that sign, hence below \(m\).
If \(n=\cdot\), the target marking is also \(\cdot\), so every recomputed
mark is below it. The result body has type below \(\rho\) by induction.
Thus the complete abstraction has type below \(\tau\).
In a neutral spine each argument remains below its receiving type;
the applications remain well typed. This proves assertion 1 and
termination of the decorated finite expansion.
Every construction step is beta conversion, general eta expansion, or
conversion congruence in a well-typed context. Thus the elaborated term
is also convertible to the original one, as asserted.

For semantic preservation, beta contractions preserve promoted values
by the substitution lemma, (C), and composition of coercions.
At the expected arrow choose \(u\) fresh for the term.
For \(a\in D_\sigma\), the promoted original function has value
\[
 \pi_{\epsilon,\rho}
       \bigl(\sem M_\varphi(\pi_{\sigma,\delta}(a))\bigr)
\]
by (C). This is the promoted value of \(M(u)\) at
\(\varphi_u^a\). By induction it is the promoted value of its elaborated
body. The abstraction interpretation and pointwise equality therefore
give exactly the promoted value of the elaborated lambda term.
For a neutral head, induction preserves every argument value at the
receiving type; successive applications then preserve the result value.
This proves assertion 2 in arbitrary coherent structures.

Finally the erased ordinary long form has a uniquely determined marked
decoration at a fixed expected type and fixed typed context.
At each lambda its binder must have the expected arrow's argument type.
At a neutral head its declaration specifies each argument's expected type.
The intrinsic abstraction marking is recomputed from the decorated body.
Induction on this finite long tree determines every annotation, modulo
fresh typed bound names. If \(M\approx N\), their erased terms have the
same ordinary long form. The decoration just described is therefore the
same for both, proving assertion 3.
Apply assertion 2 to both terms and identify their long values by typed
alpha conversion. Their promoted original values agree at the common
comparison type, which is (E).
\end{proof}
\end{annexitem}

\begin{annexitem}{annex:p7.monotonicity.term-model}
 {thm:p07-monotonicity-term-model}{Indexed functional realization and interpretation}
\begin{proof}
A class in \(T_\tau\) must always be used with a representative of intrinsic
type at most \(\tau\). If
\(\langle M\rangle\in T_{\sigma\xrightarrow m\tau}\), choose
\(M:\delta\xrightarrow n\epsilon\), where
\(\sigma\preceq\delta\), \(\epsilon\preceq\tau\), \(n\sqsubseteq m\).
Any representative \(O\) for a class of \(T_\sigma\) has
\(\operatorname{tp}(O)\preceq\sigma\), so \(M(O)\) is well typed.
Changing either representative preserves the application class by
conversion congruence; the mixed applications also type-check.
Term-domain inclusions preserve order by the preceding lemma.

Construct \(D_\tau\) and \(\iota_\tau\) by type height.
At a base type use a tagged copy of \(T_b\) with the transported preorder.
Once \(D_\sigma,D_\tau\) and their bijections are available, associate
to \(a=\langle M\rangle\in T_{\sigma\xrightarrow m\tau}\) the actual
function
\[
 F_a^{\sigma,\tau}(\iota_\sigma\langle O\rangle)
     =\iota_\tau\langle M(O)\rangle.
\]
Surjectivity of \(\iota_\sigma\) and the well-defined class action make
this a function on all of \(D_\sigma\).
Take the arrow carrier to be the set of these functions.

This association is injective. If two representing functions are equal,
evaluate them at \(\iota_\sigma\langle u\rangle\), with \(u:\sigma\)
fresh for their representatives \(M,N\).
Smaller-type injectivity gives \(M(u)\approx N(u)\).
Abstraction congruence and the original general eta rule give
\(M\approx N\). The argument includes representatives whose intrinsic
argument domains are strictly above \(\sigma\).
It proves equality of conversion classes, not uniqueness of raw marked
normal terms. Surjectivity is the definition of the arrow carrier.

Give the arrow carrier pointwise order.
Its comparison, using the smaller order isomorphisms and surjectivity at
the argument type, is exactly the term-structure extensional comparison.
Thus the bijection is an order isomorphism.
The function \(F_a^{\sigma,\tau}\) depends on its domain and codomain,
but not on \(m\). Inclusion of signed term-class domains in the unmarked
one therefore gives literal inclusion of the same actual functions.
Their required variance follows from the signed term actions.

Set \(\pi_{\alpha,\beta}=\iota_\beta i_{\alpha,\beta}\iota_\alpha^{-1}\).
Identity, composition and order preservation follow from the corresponding
properties of the inclusions. To verify (C), let
\(\alpha=\delta\xrightarrow n\epsilon\preceq
 \beta=\sigma\xrightarrow m\tau\),
\(f=\iota_\alpha\langle M\rangle\), and
\(a=\iota_\sigma\langle O\rangle\).
The left side of (C) is
\(\iota_\tau\langle M(O)\rangle\).
On the right, the input coercion merely views \(\langle O\rangle\)
in \(T_\delta\); application gives
\(\iota_\epsilon\langle M(O)\rangle\), and result coercion gives
the same \(\iota_\tau\)-value. Surjectivity at \(\sigma\) proves (C)
on every argument. This establishes every prestructure condition.

For a valuation \(\varphi\), choose for each \(x:\sigma\) a representative
\(s(x)\) of the inverse class \(\iota_\sigma^{-1}(\varphi(x))\)
with type at most \(\sigma\). Choice permits the simultaneous selection.
Define
\(\sem M_\varphi=\iota_{\operatorname{tp}(M)}\langle M[s]\rangle\).
If another lift \(t\) is chosen, injectivity gives
\(s(x)\approx t(x)\) at each variable.
Only finitely many images occur in a given term.
Repeated congruence, with fresh binders, yields \(M[s]\approx M[t]\);
all mixed substitutions are well typed by refinement.
Thus interpretation is independent of the lift.

A variable has its assigned value. A constant has the fixed value
\(\iota_{\operatorname{tp}(c)}\langle c\rangle\).
For application \(M(N)\), retain its original head type
\(\delta\xrightarrow m\epsilon\) and argument type
\(\rho\preceq\delta\). The recursive interpretation is
\[
 \iota_{\delta\xrightarrow m\epsilon}\langle M[s]\rangle
   \bigl(\pi_{\rho,\delta}
             (\iota_\rho\langle N[s]\rangle)\bigr)
 =\iota_\epsilon\langle M[s](N[s])\rangle,
\]
exactly the stipulated value of the application.
This calculation uses the original comparison indices, not possibly
refined substituted intrinsic types.

For \(\lambda x:\sigma.M\), choose a legal representative \(A\) for
an input \(a=\iota_\sigma\langle A\rangle\), and fresh \(y:\sigma\).
Application of its representing function, capture-avoiding substitution
and beta conversion give
\[
 \sem{\lambda x.M}_\varphi(a)
 =\iota_{\operatorname{tp}(M)}
      \langle(\lambda y.M[s_x^y])(A)\rangle
 =\iota_{\operatorname{tp}(M)}\langle M[s_x^A]\rangle
 =\sem M_{\varphi_x^a}.
\]
The substitution \(s_x^A\) represents precisely the changed valuation.
Its result belongs to the correct signed carrier by type refinement and
the construction of that carrier. All interpretation clauses and closure
under abstraction have now been verified.
\end{proof}
\end{annexitem}

\begin{annexitem}{annex:p7.monotonicity.completeness}
 {thm:p07-monotonicity-completeness}{The canonical model under all assignments}
\begin{proof}
For legal representatives at \(\tau\), define
\[
 \langle M\rangle\sqsubseteq_\tau\langle N\rangle
 \quad\Longleftrightarrow\quad
 \Gamma\vdash_{\mathrm{mon}}M\precsim_\tau N.
\]
Equiv in both directions followed by Trans makes the relation independent
of the representatives. Ref and Trans make it a preorder.

For a class below \(\sigma\xrightarrow+\tau\), choose a representative
of type \(\delta\xrightarrow+\epsilon\), with
\(\sigma\preceq\delta\), \(\epsilon\preceq\tau\).
Promote an input comparison to \(\delta\) by Pres, apply Mono,
and promote the resulting comparison from \(\epsilon\) to \(\tau\).
This gives signed positive action. Anti gives the negative action.
Point gives one direction of functional extensionality.
If the application comparison holds for every class in \(T_\sigma\),
it holds at a fresh variable \(u:\sigma\). Func and general eta,
as proved in the extensional-comparison lemma, give the function
comparison. Thus all term-structure conditions hold.

Realize this term structure by the preceding lemma.
For every derivable assertion \(G\precsim_\tau H\) and arbitrary
valuation \(\varphi\), choose its substitution lift \(s\).
The instantiation lemma gives
\(\Gamma\vdash G[s]\precsim_\tau H[s]\), with the same premises.
The term preorder and its order isomorphism give
\[
 \iota_\tau\langle G[s]\rangle
       \leq_\tau\iota_\tau\langle H[s]\rangle.
\]
By the definitions of coercion and interpretation, these are the promoted
values of \(G,H\) at \(\varphi\).
In particular the canonical structure satisfies every member of
\(\Gamma\) under every assignment, as the consequence definition demands.

Suppose now that the target assertion is a semantic consequence.
It holds in this canonical coherent structure at the identity-generated
valuation
\(\varphi_0(x)=\iota_{\operatorname{tp}(x)}\langle x\rangle\).
The identity substitution lifts \(\varphi_0\).
Reflection of order gives
\(\langle M\rangle\sqsubseteq_\tau\langle N\rangle\), and the
definition of this preorder gives the required derivation.
The construction uses actual functions but need not fill the arrow
carriers with all possible functions. This is why it proves the stated
general-model completeness and makes no full-hierarchy claim.
\end{proof}
\end{annexitem}

\begin{annexitem}{annex:p7.monotonicity.lyndon}
 {thm:p07-monotonicity-lyndon}{The comparison-indexed occurrence model}
\begin{proof}
Fix three distinct variables \(x,y,z:\sigma_0\).
Always choose bound variables outside this triple; bound occurrences may
equivalently be recorded by de Bruijn indices.
For a typed term, erase its colors by replacing its \emph{free}
\(y,z\) occurrences by \(x\). This is same-type free-variable renaming,
so it preserves the intrinsic type, the constructor tree, all binder
types and every corresponding occurrence polarity.

For two terms \(P,Q\) with the same color-erased typed tree \(S\),
compare the colors at each free \(x\)-position of \(S\).
At a positive position allow equality and the change \(y\) to \(z\).
At a negative position allow equality and the change \(z\) to \(y\).
At an unspecified position allow only equality.
The color \(x\) is comparable only with itself at every position.
Write \(P\leq_{\mathrm{col}}Q\) when all these conditions hold.
All other names, constants, constructors and annotations must agree.
For legal representatives at each comparison type define
\[
 \langle P\rangle\sqsubseteq_\tau\langle Q\rangle
 \quad\Longleftrightarrow\quad
 L_\tau(P)\leq_{\mathrm{col}}L_\tau(Q).
 \tag{O}
\]
The uniqueness of the comparison-directed long form makes (O)
independent of representatives.
Reflexivity is pointwise color equality. If \(P,Q\) share an erased tree
and \(Q,R\) share one, the two trees are the same. Each of the three
position-wise color orders is transitive, so \(P,R\) satisfy it.
Thus every \(T_\tau\) is preordered.

We establish three syntactic preservation facts before checking the
term-structure axioms.

First, if \(P\leq_{\mathrm{col}}Q\) and both terms can be elaborated
at \(\tau\), then
\(L_\tau(P)\leq_{\mathrm{col}}L_\tau(Q)\).
Run the same deterministic elaboration on both color-erased trees.
The three distinguished variables have identical types and none is a
binder. Their names therefore affect neither redex locations, binder
occurrence positions, nor expected types along a neutral spine.
Both sides contract or expand matching nodes, with fresh bound names
chosen for their union.
A changed color position is signed. In a beta contraction, its residual
in the abstraction body persists unless it lies in a substituted argument.
In the latter case its original signed total polarity requires a signed,
nonvacuous abstraction binder; substitution produces copies with the same
product sign. Refinement of a sign leaves that sign unchanged.
The same calculation in a surrounding context preserves a signed total
polarity. Eta expansion keeps previous free occurrences in the head and
adds only fresh bound occurrences. Processing a smaller expected binder
type uses the already proved type-refining substitution case.
Thus every surviving old color difference keeps its permitted direction,
and both color-erased outputs still agree. Any erased positions create
no new difference. Induction over the finite elaboration proves the fact.

Second, a \emph{common} typed substitution for a variable \(u\) outside
\(x,y,z\) preserves color order, and so does subsequent common-type
elaboration. In \(P\leq_{\mathrm{col}}Q\), the free \(u\)-positions,
types and polarities agree. Inserting the same term \(A\) at those
positions inserts identical copies on the two sides.
In particular, \(A\) may itself contain free \(x,y,z\); every color it
introduces is initially equal in the corresponding two copies.
An old changed distinguished occurrence is not \(u\) and retains its
signed direction under type refinement.
Fresh renaming prevents capture. The substituted color-erased trees
still agree, so the first fact applies to their subsequent elaborations.
Matching reductions duplicate or erase matching inserted copies; they
cannot turn these initially equal inserted colors into an old unequal
pair. This proves the common-substitution claim including colored
arguments.

Third, suppose all free occurrences of \(u\) in one term \(B\) have
sign \(\epsilon\in\{+,-\}\), and \(N\leq_{\mathrm{col}}O\), with both
legal at \(u\)'s type.
If \(\epsilon=+\), then
\[
 B[N/u]\leq_{\mathrm{col}}B[O/u];
\]
if \(\epsilon=-\), then
\[
 B[O/u]\leq_{\mathrm{col}}B[N/u].
\]
To verify this, the two substituends have the same erased typed tree and
the same intrinsic type. Hence the two resulting substitutions have the
same erased typed tree. Distinguished colors already in \(B\) are equal
on both sides. A changed substituend occurrence of sign \(\delta\)
has total sign \(\epsilon\circ\delta\) at each copy, by the signed
instance of the substitution lemma.
For \(\epsilon=+\), its allowed direction is unchanged.
For \(\epsilon=-\), its sign reverses, exactly compensated by comparing
the two substituted terms in the reverse order.
These are every possible changed position. The first fact gives the
same comparisons after elaboration. These three facts concern syntax
and signed residuals; they use no semantic converse.

We now check signed action. Let
\(\langle M\rangle\in T_{\sigma\xrightarrow+\tau}\) and
\(\langle N\rangle\sqsubseteq_\sigma\langle O\rangle\).
Write, using one fresh binder,
\[
 L_{\sigma\xrightarrow+\tau}(M)=\lambda u:\sigma.B,\qquad
 N'=L_\sigma(N),\quad O'=L_\sigma(O).
\]
Type validity of the long form forces this abstraction's intrinsic
marking to be \(+\). Its binder is therefore nonvacuous and uniformly
positive. Definition (O) gives \(N'\leq_{\mathrm{col}}O'\).
The third preservation fact and the first give the comparison of the
elaborated \(B[N'/u]\) and \(B[O'/u]\) at \(\tau\).
The original conversion relation relates a term to its long elaboration;
application congruence and beta conversion identify those outputs with
\(L_\tau(M(N))\) and \(L_\tau(M(O))\).
Thus positive term action is monotone.
At a negative arrow the long binder is uniformly negative, so the same
argument reverses the output comparison. This proves negative action.
No prior order-preservation assumption about type inclusions was used.

For functional extensionality choose the long forms of two classes in
\(T_{\sigma\xrightarrow m\tau}\) with the same fresh binder:
\[
 L_{\sigma\xrightarrow m\tau}(M)=\lambda u:\sigma.B,\qquad
 L_{\sigma\xrightarrow m\tau}(N)=\lambda u:\sigma.C.
\]
The bodies are respectively \(L_\tau(M(u))\) and \(L_\tau(N(u))\).
Deleting or adding this common lambda does not change the polarity
of any distinguished free occurrence.
If the arrow classes are ordered, \(B\leq_{\mathrm{col}}C\).
For any legal argument class, substitute its long representative for
\(u\) on both sides. The second preservation fact, including its
colored-argument case, and conversion identify the normalized outputs
with the desired application comparison at \(\tau\).
Conversely, if the application comparison holds at every argument class,
it holds at the fresh variable \(u:\sigma\).
Its defining comparison is \(B\leq_{\mathrm{col}}C\).
Adding the common binder gives the arrow comparison.
This proves both directions of functional extensionality.

We have established all preorder and signed-action axioms of a term
structure. The term-inclusion lemma now supplies order preservation
between its types. Realize it by
Lemma~\ref{thm:p07-monotonicity-term-model}.
This is a coherent structure for the full marked type system.

In this structure
\(\langle y\rangle\sqsubseteq_{\sigma_0}\langle z\rangle\).
At a base type this is the permitted positive color change.
At an arrow type, long expansion adds fresh bound arguments to the
distinguished free head, which retains positive polarity.
Iterating through its type gives the same permitted change.
The realization therefore orders
\(\iota_{\sigma_0}\langle y\rangle\leq
  \iota_{\sigma_0}\langle z\rangle\).

Suppose the promoted interpretation of \(M\) is monotone in \(x\)
in every coherent structure. Choose \(y,z\) fresh for \(M\), use this
constructed structure, and use the identity-generated assignment on
all other variables. Vary \(x\) between the two ordered values just given.
The indexed interpretation and output promotion yield
\[
 \iota_\tau\langle M[y/x]\rangle
       \leq_\tau
 \iota_\tau\langle M[z/x]\rangle.
\]
Order reflection and (O) turn this into
\[
 L_\tau(M)[y/x]\leq_{\mathrm{col}}L_\tau(M)[z/x].
 \tag{*}
\]
Same-type free renaming commutes with elaboration: it changes neither
types, redex positions, expected-type choices nor binding structure,
and fresh bound variables are chosen uniformly.
This justifies the two displayed renamed long forms.
After color erasure their common tree is \(L_\tau(M)\).
At every free \(x\)-position, (*) changes \(y\) into \(z\);
the color order permits that distinct change only at a positive position.
Thus every such occurrence is positive.
Under universal antitonicity the output comparison reverses; every
changed position goes from \(z\) to \(y\), which is permitted only
at a negative position. This proves both converses, also when the
set of positions is empty.

For the forward directions use semantic preservation of \(L_\tau(M)\),
apply the proved occurrence-polarity implication to that elaborated term,
and, if necessary, compose its value with the order-preserving output
coercion to \(\tau\). The resulting map is exactly the promoted map
in the theorem. An absent variable gives an assignment-independent
value by coincidence, hence a constant map. This proves every stated case.
\end{proof}
\end{annexitem}
